% ═══════════════════════════════════════════════════════════════
% MT-05b: MINITEST VERBEN -ER/-IR
% Spanisch Klasse 7 - Sequenz 5b
% ═══════════════════════════════════════════════════════════════
% Dauer: 15 Minuten
% Punkte: 28 BE
% ═══════════════════════════════════════════════════════════════

\documentclass[11pt, a4paper]{article}

\newif\ifswmodus
\swmodusfalse

\usepackage[utf8]{inputenc}
\usepackage[T1]{fontenc}
\usepackage[ngerman]{babel}
\usepackage{helvet}
\renewcommand{\familydefault}{\sfdefault}

\usepackage[margin=2cm]{geometry}
\usepackage{enumitem}
\usepackage{tabularx}
\usepackage{booktabs}

\usepackage{xcolor}
\usepackage{tcolorbox}
\tcbuselibrary{skins}

\usepackage{fancyhdr}
\usepackage{lastpage}
\usepackage{needspace}

% FARBEN
\definecolor{spanisch-color}{HTML}{c41e3a}
\definecolor{grau-color}{HTML}{888888}
\definecolor{hellgrau-color}{HTML}{f5f5f5}
\definecolor{er-color}{HTML}{2c5aa0}
\definecolor{ir-color}{HTML}{2a7a4b}

\definecolor{sw-dark}{HTML}{000000}
\definecolor{sw-gray}{HTML}{666666}
\definecolor{sw-white}{HTML}{ffffff}

\ifswmodus
    \colorlet{spanisch}{sw-dark}
    \colorlet{grau}{sw-gray}
    \colorlet{hellgrau}{sw-white}
    \colorlet{erfarbe}{sw-dark}
    \colorlet{irfarbe}{sw-dark}
\else
    \colorlet{spanisch}{spanisch-color}
    \colorlet{grau}{grau-color}
    \colorlet{hellgrau}{hellgrau-color}
    \colorlet{erfarbe}{er-color}
    \colorlet{irfarbe}{ir-color}
\fi

\newcommand{\fachfarbe}{spanisch}

\newcommand{\thema}{Minitest: Verben -er/-ir}
\newcommand{\fach}{Spanisch}
\newcommand{\klasse}{7}
\newcommand{\maxpunkte}{28}
\newcommand{\zeit}{15 Minuten}
\newcommand{\datum}{\today}

\pagestyle{fancy}
\fancyhf{}
\renewcommand{\headrulewidth}{0.4pt}
\renewcommand{\footrulewidth}{0.4pt}

\lhead{\fach{} -- Klasse \klasse}
\chead{\textbf{MT-05b}}
\rhead{\datum}

\lfoot{\zeit}
\cfoot{}
\rfoot{Seite \thepage\ von \pageref{LastPage}}

\newcommand{\punktebox}[1]{\hfill\fbox{\small \_/#1}}
\newcommand{\luecke}[1]{\rule{#1}{0.4pt}}

\newtcolorbox{infobox}{
    colback=hellgrau,
    colframe=\fachfarbe,
    boxrule=1pt,
    arc=0pt,
    left=8pt, right=8pt, top=8pt, bottom=8pt
}

\newenvironment{aufgabe}[1]{%
    \needspace{5\baselineskip}%
    \nopagebreak[4]%
    \vspace{10pt}
    \noindent\textbf{\textcolor{\fachfarbe}{Aufgabe #1}}
    \nopagebreak[4]%
    \vspace{5pt}
    \nopagebreak[4]%
    \begin{list}{}{\leftmargin=0pt}
    \item[]
}{%
    \end{list}
}

\newcommand{\es}[1]{\textcolor{\fachfarbe}{\textbf{#1}}}
\newcommand{\er}[1]{\textcolor{erfarbe}{\textbf{#1}}}
\newcommand{\ir}[1]{\textcolor{irfarbe}{\textbf{#1}}}

\begin{document}

% ═══════════════════════════════════════════════════════════════
% KOPFBEREICH
% ═══════════════════════════════════════════════════════════════

\begin{center}
    {\Large\bfseries\textcolor{\fachfarbe}{\thema}}
\end{center}

\vspace{5pt}

\noindent
\begin{tabularx}{\textwidth}{|X|X|X|}
\hline
\textbf{Name:} \luecke{4cm} &
\textbf{Klasse:} \klasse &
\textbf{Datum:} \luecke{2cm} \\
\hline
\end{tabularx}

\vspace{10pt}

\begin{infobox}
\textbf{Zeit:} \zeit{} | \textbf{Punkte:} \maxpunkte{} BE | \textbf{Hilfsmittel:} keine
\end{infobox}

\vspace{15pt}

% ═══════════════════════════════════════════════════════════════
% AUFGABE 1: -er Verben konjugieren (8P)
% ═══════════════════════════════════════════════════════════════

\begin{aufgabe}{1 -- Verben auf \er{-er} konjugieren}
Setze die richtige Verbform ein. \punktebox{8}
\end{aufgabe}

\noindent
a) Yo \luecke{3cm} una pizza. \er{(comer)}\\[1em]
b) Tú \luecke{3cm} agua. \er{(beber)}\\[1em]
c) María \luecke{3cm} un libro. \er{(leer)}\\[1em]
d) Nosotros \luecke{3cm} en el parque. \er{(correr)}

\vspace{15pt}

% ═══════════════════════════════════════════════════════════════
% AUFGABE 2: -ir Verben konjugieren (8P)
% ═══════════════════════════════════════════════════════════════

\begin{aufgabe}{2 -- Verben auf \ir{-ir} konjugieren}
Setze die richtige Verbform ein. \punktebox{8}
\end{aufgabe}

\noindent
a) Yo \luecke{3cm} en Rostock. \ir{(vivir)}\\[1em]
b) Tú \luecke{3cm} un mensaje. \ir{(escribir)}\\[1em]
c) Pedro \luecke{3cm} la ventana. \ir{(abrir)}\\[1em]
d) Nosotros \luecke{3cm} en Alemania. \ir{(vivir)}

\vspace{15pt}

% ═══════════════════════════════════════════════════════════════
% AUFGABE 3: Endungen ergänzen (6P)
% ═══════════════════════════════════════════════════════════════

\begin{aufgabe}{3 -- Endungen ergänzen}
Ergänze die fehlenden Endungen. \punktebox{6}
\end{aufgabe}

\begin{center}
\begin{tabular}{|l|c|c|}
\hline
\textbf{Person} & \er{-er (comer)} & \ir{-ir (vivir)} \\
\hline
yo & com\luecke{1.5cm} & viv\luecke{1.5cm} \\[0.8em]
\hline
nosotros & com\luecke{1.5cm} & viv\luecke{1.5cm} \\[0.8em]
\hline
ellos & com\luecke{1.5cm} & viv\luecke{1.5cm} \\[0.8em]
\hline
\end{tabular}
\end{center}

\vspace{15pt}

% ═══════════════════════════════════════════════════════════════
% AUFGABE 4: Sätze bilden (6P)
% ═══════════════════════════════════════════════════════════════

\begin{aufgabe}{4 -- Sätze bilden}
Bilde Sätze. Konjugiere das Verb! \punktebox{6}
\end{aufgabe}

\noindent
a) yo / leer / un libro

\luecke{10cm}

\vspace{1em}

\noindent
b) mi hermano / correr / en el parque

\luecke{10cm}

\vspace{1em}

\noindent
c) nosotros / vivir / en Alemania

\luecke{10cm}

% ═══════════════════════════════════════════════════════════════
% PUNKTEÜBERSICHT
% ═══════════════════════════════════════════════════════════════

\vspace{25pt}
\hrule
\vspace{10pt}

\noindent
\begin{tabularx}{\textwidth}{|X|c|c|c|c||c|}
\hline
\textbf{Aufgabe} & 1 & 2 & 3 & 4 & \textbf{Gesamt} \\
\hline
\textbf{max. Punkte} & 8 & 8 & 6 & 6 & \textbf{28} \\
\hline
\textbf{erreicht} & & & & & \\[0.8em]
\hline
\end{tabularx}

\vspace{10pt}

\begin{flushright}
\textbf{Note:} \luecke{2cm}
\end{flushright}

\end{document}
